\subsection{万有覆叠}

定义 1. —— 带基点集是指带有其一个元素的集合 X，该元素称为基点。带有元素 x 的集合 X 有时记作 (X, x)。设 (X, x) 与 (Y, y) 为带基点集；从 (X, x) 到 (Y, y) 的带基点映射是指满足 \(f(x) = y\) 的从 X 到 Y 的映射 f。

带基点拓扑空间、带基点连续映射、带基点拓扑空间的带基点覆叠等概念，都可类似地定义。

若 \((X,x)\) 与 \((Y,y)\) 为带基点拓扑空间，则从 \((X,x)\) 到 \((Y,y)\) 的带基点连续映射组成的集合记作 \(\mathcal{C}((X,x);(Y,y))\)。

人们常不用"设 f 为从 \((X, x)\) 到 \((Y, y)\) 上的带基点映射"这种说法，而使用下述短语："设 \(f: (X, x) \to (Y, y)\) 为带基点映射"。

设 B 为拓扑空间，b 为 B 的一点。

定义 2. —— 称 (B,b) 的带基点覆叠 (E,x) 为万有覆叠，如果对 (B,b) 的每个带基点覆叠 (E',x')，存在唯一的 B 的覆叠态射 \(f: \mathrm{E} \to \mathrm{E}'\) ，使得 \(f(x) = x'\) 。

若 (E, x) 与 (E', x') 都是 (B, b) 的万有覆叠，则从 (E, x) 到 (E', x') 的唯一的 B-态射是 B-空间的同构。

设 E 为 B 的连通覆叠，x 为纤维 \(E_{b}\) 的一点。假设对每个带基点覆叠 \((\mathrm{E}^{\prime}, x^{\prime})\) （\((\mathrm{B}, b)\) 的），都存在 B 的覆叠态射 \(f: E \to E^{\prime}\) 使得 \(f(x) = x^{\prime}\) 。这样的态射 f 于是是唯一的（I, 第 34 页, 命题 11 的推论 1），故 \((\mathrm{E}, x)\) 是 \((\mathrm{B}, b)\) 的万有覆叠。*我们以后将会看到（I, 第 126 页, 命题 3 的推论），特别地，当 E 的一切覆叠都可平凡化时即属此情形。*

命题 1. —— 设 B 为连通且局部连通的拓扑空间，b 为 B 的一点。设 (E, x) 为 (B, b) 的万有覆叠。则 E 是 B 的伽罗瓦覆叠，且 B 的每个覆叠都与 E 相伴。

空间 E 是局部连通的，因为 B 是局部连通的。设 \(E_{0}\) 为 x 在 E 中的连通分支，于是空间 \((\mathrm{E}_{0}, x)\) 是 \((\mathrm{B}, b)\) 的带基点覆叠（I, 第 80 页, 命题 6 的推论 1）。于是存在唯一的覆叠态射 \(f: E \to E_{0}\) 使得 \(f(x) = x\)。若以 i 表示 \(E_{0}\) 到 E 中的典范单射，则映射 \(i \circ f: E \to E\) 是把 x 映到 x 的覆叠态射，映射 \(Id_{E}\) 亦然；由于 \((\mathrm{E}, x)\) 是 \((\mathrm{B}, b)\) 的万有覆叠，我们有 \(i \circ f = Id_{E}\)。这蕴含 i 是满射，从而 \(E_{0} = E\)。于是 E 是连通的。

设 \(y\) 为 \(\mathrm{E}_b\) 的一点，并考察带基点覆叠 \((\mathrm{E},y)\) （\((\mathrm{B},b)\) 的）；依假设存在唯一的覆叠态射 \(f\colon \mathrm{E}\to \mathrm{E}\) 使得 \(f(x) = y\) 。映射 \(s\colon \mathrm{E}\to \mathrm{E}\times_{\mathrm{B}}\mathrm{E}\) 由 \(t\mapsto (t,f(t))\) 所定义，它是映射 \(\mathrm{pr}_1\colon \mathrm{E}\times_{\mathrm{B}}\mathrm{E}\to \mathrm{E}\) 的连续截面。由 I, 第 81 页推论 4，这表明覆叠 \(\mathrm{E}\times_{\mathrm{B}}\mathrm{E}\) （由映射 \(\mathrm{pr}_1\) 给出的 E 上的覆叠）可平凡化。于是覆叠 E 是伽罗瓦的（I, 第 102 页, 定理 2）。再由 I, 第 112 页命题 10 的推论，B 的每个覆叠都与 E 相伴。

推论. —— 设 B 为局部连通拓扑空间。设 b 为 B 的一点，(E, x) 为 (B, b) 的万有覆叠。对 B 的子空间 A，下面两条性质等价：

\begin{enumerate}
\def\labelenumi{(\roman{enumi})}
\item
  覆叠 E 在 A 上可平凡化；
\item
  B 的每个覆叠都在 A 上可平凡化。
\end{enumerate}

性质 (ii) 显然蕴含性质 (i)。其逆由"B 的每个覆叠都与覆叠 E 相伴"这一事实得出（I, 第 105 页, 命题 7）。

